% Evans 12.6, printed pp.647--651; English PDF pp.664--668.
% Chinese start: printed p.459 / scan p.477.
\section{Problems}
\label{sec:12-6}
These are abbreviated descriptions of all 18 source problems, followed
by reconstructed solutions. Regularity, decay, and parameter
qualifications are included where the source wording suppresses them.

\subsection*{1. Energy for a gradient-dependent principal part}
For $u_{tt}-\operatorname{div}DL(Du)=0$, use
\[
 E_L(t)=\int\{u_t^2/2+L(Du)-L(0)\}dx.
\]
The subtraction makes the integral finite for compactly supported $u$.
Its derivative is $\int u_tu_{tt}+DL(Du)\cdot Du_t
=\int u_t\{u_{tt}-\operatorname{div}DL(Du)\}=0$.
This needs no convexity of $L$ for the identity; positivity or
hyperbolicity would require further assumptions.

\subsection*{2. Klein--Gordon energy and equipartition}
For $\PDEWave u+m^2u=0$, take $m\ge0$ and
$E=\frac12\int(u_t^2+|Du|^2+m^2u^2)$. Multiplication by $u_t$
proves conservation. For finite-energy data let
$\omega(\xi)=\sqrt{|\xi|^2+m^2}$ and use a unitary Fourier transform:
\[
 \widehat u=\widehat g\cos(\omega t)
                   +\widehat h\sin(\omega t)/\omega.
\]
The difference between potential and kinetic terms is
\[
 \begin{split}
 D(t)&=\int(\omega^2|\widehat u|^2-|\widehat u_t|^2)d\xi\\
 &=\int(\omega^2|\widehat g|^2-|\widehat h|^2)\cos(2\omega t)
       +2\omega\operatorname{Re}(\widehat g\overline{\widehat h})
                                      \sin(2\omega t)\,d\xi.
 \end{split}
\]
Both amplitudes are $L^1$ by Cauchy--Schwarz and finite energy.
In polar coordinates change the radial variable from $r$ to
$\omega=\sqrt{r^2+m^2}$. The resulting amplitudes are $L^1$ on
$(m,\infty)$; the possible critical point $r=0$ has no atom.
Riemann--Lebesgue gives $D(t)\to0$. Since potential plus kinetic
is $2E$, each tends to $E$, including the claimed
$\int(|Du|^2+m^2u^2)\to E$. A finite-energy approximation removes
any auxiliary smoothness used for the Fourier calculation.

\subsection*{3. Descent and the one-dimensional massive kernel}
Put $\overline u(x,y,t)=u(x,t)\cos(my)$. Then
$(\partial_t^2-\Delta_x-\partial_y^2)\overline u
=\cos(my)(\PDEWave u+m^2u)=0$.
For $n=1$, apply the two-dimensional Poisson formula and integrate
out the additional variable. With
$J_0(z)=\pi^{-1}\int_0^\pi\cos(z\sin\theta)d\theta$, the result is
\[
 \begin{split}
 u(x,t)={}&\frac12\partial_t\int_{x-t}^{x+t}
       g(y)J_0\!\left(m\sqrt{t^2-(x-y)^2}\right)dy\\
 &+\frac12\int_{x-t}^{x+t}
       h(y)J_0\!\left(m\sqrt{t^2-(x-y)^2}\right)dy.
 \end{split}
\]
The integrated kernel identity is
$\int_{-a}^a\cos(my)(a^2-y^2)^{-1/2}dy=\pi J_0(ma)$,
obtained by $y=a\sin\theta$.
If desired, differentiating the first integral gives endpoint term
$\{g(x+t)+g(x-t)\}/2$ and interior kernel
$-mtJ_1(m\sqrt{t^2-(x-y)^2})/(2\sqrt{t^2-(x-y)^2})$.
Its endpoint limit is finite because $J_1(z)\sim z/2$.
At $m=0$ this reduces exactly to d'Alembert's formula.

\subsection*{4. Removing constant damping}
For $\PDEWave u+du_t=0$, put $v=e^{dt/2}u$.
Direct substitution gives $\PDEWave v-(d^2/4)v=0$.
Thus the constant zero-order coefficient is negative when $d>0$;
the damping has not turned into a positive Klein--Gordon mass.

\subsection*{5. Phase speed and the domain of dependence}
The plane wave $e^{i(x\cdot\xi-\omega t)}$ has dispersion
$\omega^2=|\xi|^2+m^2$. For $m>0$ its phase speed is
$\omega/|\xi|>1$, but its group velocity is
$D_\xi\omega=\xi/\omega$, of magnitude less than one.
A plane wave has no compact support and carries no newly created
front. Causal propagation of a localized disturbance is the cone
statement of \S\ref{sec:12-1}, which remains true for the massive
lower-order term. Group velocity describes narrow-band packets;
the finite propagation proof is what controls an actual front.
For $m=0$ both speeds have magnitude one.

\subsection*{6. Uniform three-dimensional decay}
Kirchhoff gives
\[
 u(x,t)=\frac1{4\pi}\int_{S^2}
 [g(x+t\omega)+tDg(x+t\omega)\cdot\omega+th(x+t\omega)]d\omega.
\]
If the data are supported in $B(0,R)$ and $t\ge2R$, the part of
$S^2$ where $x+t\omega\in B(0,R)$ has area at most $C R^2/t^2$.
This follows by intersecting the sphere of radius $t$ centered at
$x$ with that fixed ball and using the spherical-cap area formula;
if there is no intersection the integral is zero.
The displayed formula is therefore bounded by $C/t$ uniformly in
$x$. For $t\le2R$, the same formula with the full sphere bounds
$u$ by a fixed constant, which is also at most a larger $C/t$.

\subsection*{7. Uniform two-dimensional decay}
Differentiate the rescaled Poisson formula to write
\[
 u(x,t)=\frac1{2\pi}\int_{|x-y|<t}
 \frac{h(y)+\{g(y)+(y-x)\cdot Dg(y)\}/t}
                       {\sqrt{t^2-|x-y|^2}}dy.
\]
For $t$ larger than a fixed multiple of the support radius, its
absolute value is bounded by a data constant times
$J=\int_{B_R\cap B(x,t)}(t^2-|x-y|^2)^{-1/2}dy$.
On $|x-y|\le t/2$ the contribution is $C_R/t$.
For the remaining radii use coarea about $x$. The length of a circle
of radius $r>t/2$ intersected with $B_R$ is at most $C R$;
the possible radii form an interval of length at most $2R$.
Since the kernel increases toward $r=t$, this contribution is at most
\[
 C R\int_{\max(t-2R,t/2)}^t\frac{dr}{\sqrt{t^2-r^2}}
 \le C_Rt^{-1/2}.
\]
The bounded small-time strip follows directly from the rescaled
Poisson formula and the smooth data. Hence $|u(x,t)|\le Ct^{-1/2}$
for all $x,t>0$. Unlike the three-dimensional case, the interior
of the light cone generally carries a tail.

\subsection*{8. Bounded displacement data can amplify}
For the given radial $g^\varepsilon$, $0\le g^\varepsilon\le1$
and $g^\varepsilon(r)\to\mathbf1_{(1,3)}(r)$ a.e. as
$\varepsilon\downarrow0$. At $x=0$ and fixed $t>3$, Poisson gives
\[
 u^\varepsilon(0,t)=-t\int_1^3
        \frac{r g^\varepsilon(r)}{(t^2-r^2)^{3/2}}dr
 \longrightarrow t\left\{\frac1{\sqrt{t^2-1}}
                         -\frac1{\sqrt{t^2-9}}\right\}.
\]
Dominated convergence is valid for each such fixed $t$.
The limit tends to $-\infty$ as $t\downarrow3$.
Given any $M$, first choose $3<t<4$ making its magnitude exceed
$2M$, then take $\varepsilon$ small enough. This proves that the
supremum in the problem tends to infinity. No maximum principle
for wave displacement follows from a bound on $g$ alone.

\subsection*{9. Hyperbolic Kelvin transform}
Let $d=n+1$, $X=(x,t)$, $q=|x|^2-t^2\ne0$,
$Y=X/q$, $a=(d-2)/2$ and $w(X)=|q|^{-a}$.
For the Lorentz matrix $G=\operatorname{diag}(-1,\ldots,-1,1)$,
direct differentiation gives
\[
 DY\,G\,(DY)^T=q^{-2}G,\qquad \PDEWave w=0,
 \qquad w\PDEWave Y^j+2\langle Dw,DY^j\rangle_G=0.
\]
For example $\PDEWave q=-2d$, $\langle Dq,Dq\rangle_G=-4q$,
$\PDEWave Y^j=2(d-2)X^j/q^2$, and
$\langle Dw,DY^j\rangle_G=-2awX^j/q^2$.
The product and chain rules therefore cancel every first-derivative
term and leave
\[
 \PDEWave\{w\,u(Y)\}=|q|^{-(d+2)/2}(\PDEWave u)(Y).
\]
It follows that the transform preserves free waves on each region
where $q\ne0$. Absolute values allow both signs of $q$;
no extension through the characteristic cone is asserted.

\subsection*{10. Compatibility of the sine--Gordon system}
Set $A=(u+v)/2$, $B=(u-v)/2$. The given equations are
$B_t=a\sin A$, $A_x=a^{-1}\sin B$.
Differentiating gives $B_{xt}=\cos A\sin B$ and
$A_{xt}=\cos B\sin A$. Adding and subtracting yield
$u_{xt}=\sin(A+B)=\sin u$ and
$v_{xt}=\sin(A-B)=\sin v$.
These $x,t$ are characteristic coordinates. With new variables
$T=x+t$, $X=x-t$, the mixed derivative becomes
$\partial_x\partial_t=\partial_T^2-\partial_X^2$.
Thus the equation has the usual wave-operator form in $(X,T)$,
with no missing factor when these summed coordinates are used.

\subsection*{11. A B\"acklund kink from the zero solution}
With $v=0$, the equations become
$u_t=2a\sin(u/2)$, $u_x=2a^{-1}\sin(u/2)$.
They imply $u=U(s)$, $s=at+x/a$, and $U'=2\sin(U/2)$.
Since $(\tan(U/4))'=\tan(U/4)$, nonconstant solutions are
\[
 u(x,t)=4\arctan(Ce^{at+x/a}),\qquad C\ne0,
\]
with compatible $4\pi$ shifts; the constant solutions $u=2\pi k$
are included separately. Differentiation gives $U''=\sin U$, so
$u_{xt}=U''=\sin u$. For $a>0,C>0$ the profile connects $0$ to
$2\pi$ as $s$ increases. The parameter $C$ translates the kink.

\subsection*{12. The Sobolev product estimate}
For $b_j=|\beta_j|$, the proof of \eqref{eq:12-product} gives
$D^{\beta_j}u_j\in L^{p_j}$ for every inverse exponent between
$b_j/(2k)$ and $1/2$, when $k>n/2$.
Choose these inverse exponents with sum exactly $1/2$; the endpoint
sums bracket $1/2$ because $\sum b_j\le k$.
Holder then gives
$\|\prod D^{\beta_j}u_j\|_2\le C\prod\|u_j\|_{H^k}$.
This handles also zero multiindices and the case of a derivative
of order exactly $k$, without using the false endpoint embedding
$H^{n/2}\hookrightarrow L^\infty$.

\subsection*{13. Sphere-valued wave maps}
Since $|u|^2=1$, differentiation gives
$u\cdot u_{tt}=-|u_t|^2$ and $u\cdot\Delta u=-|Du|^2$.
The normal space of the sphere is spanned by $u$ itself. Thus
$\PDEWave u=\gamma u$ with
$\gamma=u\cdot\PDEWave u=|Du|^2-|u_t|^2$, as asserted.
This is a genuine derivative nonlinearity, so the scalar potential
energy theory is not directly applicable.

\subsection*{14. Wave-map energy and a support correction}
A unit-length map cannot itself vanish outside a compact set.
Interpret the source's support wording as $u$ equal to a fixed unit
vector outside a compact set, or impose decay on $Du,u_t$ sufficient
for integration by parts. Then
\[
 \frac d{dt}\frac12\int(|u_t|^2+|Du|^2)
 =\int u_t\cdot\PDEWave u
 =\int(|Du|^2-|u_t|^2)u_t\cdot u=0,
\]
because $u_t\cdot u=0$. This states the nonvacuous finite-energy
version of the exercise.
\PDEAnnotation{球面值映射不能紧支为零}{若 $|u|=1$，则
$u\not\equiv0$ 于任何外部区域。题 14 的紧支应理解成
$u-u_\infty$ 紧支或导数充分衰减。能量只积分 $Du,u_t$，
不积分处处为 1 的 $|u|^2$。}

\subsection*{15. An elementary small-energy quintic proof}
Modify the average estimate in \eqref{eq:12-cone-hardy} using
$\int_{B_t}|g|\le Ct^{5/2}\|g\|_6$. Then
\[
 |\overline w|\le Ct^{-1/2}(E^{1/2}+\|g\|_6),\qquad
 \|w/r\|_2\le C(E^{1/2}+E^{1/6}),
\]
since $\|g\|_6^6\le6E$. Cone flux also gives $\int w^6\le6E$.
For $M(\tau)=\sup_{t\le\tau,x}|u(x,t)|$,
\[
 \int\frac{|w|^5}{r}
 \le\|w/r\|_2\left(\int|w|^8\right)^{1/2}
 \le C(E+E^{2/3})M(\tau).
\]
The retarded formula gives
$M(\tau)\le\|v\|_\infty+C(E+E^{2/3})M(\tau)$.
Choose a universal small initial-energy threshold making the last
coefficient at most $1/2$. On every finite strip the free wave is
bounded, so the amplitude is bounded and the smooth solution
continues globally. This argument needs neither the Morawetz identity
nor Strichartz, and makes no smallness assumption on a high Sobolev
norm of the data.

\subsection*{16. Galilean invariance for NLS}
For $iu_t+\Delta u=f(|u|^2)u$, set
$\theta=\xi\cdot x/2-|\xi|^2t/4$, $y=x-\xi t$,
$w=e^{i\theta}u(y,t)$. Then
\[
 \begin{split}
 iw_t&=e^{i\theta}(|\xi|^2u/4+iu_t-i\xi\cdot Du),\\
 \Delta w&=e^{i\theta}(\Delta u+i\xi\cdot Du-|\xi|^2u/4).
 \end{split}
\]
The extra terms cancel, and $|w|=|u(y,t)|$ makes the nonlinear
term transform with the same phase. This proves the stated invariance.

\subsection*{17. NLS mass, energy, and momentum}
Assume the solution and the needed derivatives decay rapidly.
Let $\rho=|u|^2$, $F(0)=0$, $F'=f$, with $f$ real. The mass
identity follows from
$\rho_t=-2\operatorname{div}\operatorname{Im}(\overline uDu)$.
For the energy $\mathcal H=\int(|Du|^2+F(\rho))$,
\[
 \mathcal H'=2\operatorname{Re}\int
 (-\Delta\overline u+f(\rho)\overline u)u_t
 =2\operatorname{Re}\int(-i\overline u_t)u_t=0.
\]
For $P_j=\operatorname{Im}\int\overline u\,u_{x_j}$, integration
by parts in one of the two time-derivative terms gives
\[
 P_j'=2\operatorname{Im}\int\overline u_tu_{x_j}
 =-2\operatorname{Re}\int\Delta\overline u\,u_{x_j}
       +2\operatorname{Re}\int f(\rho)\overline u\,u_{x_j}=0.
\]
The last two integrals are respectively integrals of
$\partial_j|Du|^2$ and $\partial_jF(\rho)$, hence vanish.
This agrees with the source's momentum integrand
$(\overline uDu-uD\overline u)/(2i)$.

\subsection*{18. NLS virial identity and focusing cubic blow-up}
In addition assume finite variance $V(t)=\int|x|^2|u|^2$ and
sufficient decay for the weighted integrations. The mass continuity
identity gives
$V'=4\operatorname{Im}\int x\cdot\overline uDu$.
Differentiate this expression, substitute
$u_t=i\Delta u-if(\rho)u$, and integrate twice in the Laplacian
terms. They give $2\int|Du|^2$ in the derivative of $V'/4$.
The nonlinear terms give
\[
 -\int x\cdot\rho Df(\rho)
 =-\int x\cdot D\{\rho f(\rho)-F(\rho)\}
 =n\int\{\rho f(\rho)-F(\rho)\}.
\]
Thus
\begin{equation}\label{eq:12-nls-virial}
 V''=8\int|Du|^2+4n\int\{\rho f(\rho)-F(\rho)\}.
\end{equation}
For focusing cubic NLS, $f(\rho)=-\rho$, $F=-\rho^2/2$ and
$\mathcal H=K-P/2$, with $K=\int|Du|^2$, $P=\int|u|^4$.
Hence $V''=8K-2nP=8\mathcal H+(4-2n)P\le8\mathcal H(0)<0$
when $n\ge2$ and the initial energy is negative.
Integrating gives
$0\le V(t)\le V(0)+V'(0)t+4\mathcal H(0)t^2$,
which fails at a finite time. This proves nonexistence of a global
solution in the stated regularity and finite-variance class.
The variance hypothesis is essential to this particular proof;
negative energy alone is not a license to differentiate an infinite
weighted integral.
